%% =====================================================================
\chapter{Frontend}\label{ch:frontend}

\section{Structure}

The frontend is plain HTML, CSS and JavaScript with no build step. The API
server serves it from \file{frontend\codesl{}}.

\begin{table}[H]
\centering\small
\caption{Frontend files}\label{tab:frontend}
\begin{tabularx}{\textwidth}{@{}P{4cm}Y@{}}
\toprule
\textbf{File} & \textbf{Role} \\
\midrule
\file{homepage\codedot{}html}, \file{homepagescript\codedot{}js} & Public landing page; live counts from \code{GET \codesl{}api\codesl{}public\codesl{}stats} \\
\file{login\codedot{}html}, \file{loginscript\codedot{}js} & Sign in (password, then OTP), registration for three roles \\
\file{patient\codedot{}html}, \file{doctor\codedot{}html}, \file{admin\codedot{}html}, \file{insurance\codedot{}html} & Empty shells: they load the shared scripts and one role script \\
\file{js\codesl{}api\codedot{}js} & The only code that talks to the server (tokens, refresh, errors) \\
\file{js\codesl{}ui\codedot{}js} & Escaping, formatting, loading/empty/error states, tables, modals, toasts \\
\file{js\codesl{}portal\codedot{}js} & Dashboard shell: session restore, sidebar, routing, notifications, profile \\
\file{js\codesl{}patient\codedot{}js}, \file{doctor\codedot{}js}, \file{admin\codedot{}js}, \file{insurance\codedot{}js} & One function per page of that role's dashboard \\
\file{css\codesl{}portal\codedot{}css} & Shared design system (light and dark), responsive layout \\
\bottomrule
\end{tabularx}
\end{table}

A dashboard shell is only a few lines: it loads two pinned CDN libraries
(icons and charts) and the four scripts, all with \code{defer}.

\begin{lstlisting}[style=shell,caption={Script tags of \file{frontend/patient.html}}]
<script src="https://unpkg.com/lucide@0.460.0/dist/umd/lucide.min.js" defer></script>
<script src="https://cdn.jsdelivr.net/npm/chart.js@4.4.1/dist/chart.umd.min.js" defer></script>
<script src="js/api.js" defer></script>
<script src="js/ui.js" defer></script>
<script src="js/portal.js" defer></script>
<script src="js/patient.js" defer></script>
\end{lstlisting}

\section{The API client: \file{js\codesl{}api\codedot{}js}}

\begin{lstlisting}[language=JavaScript,caption={request with transparent refresh (\file{frontend/js/api.js})}]
  async function request(method, path, { body, query, form, retry = true } = {}) {
    // ... build the URL from path + query
    const headers = {};
    if (accessToken) headers.Authorization = `Bearer ${accessToken}`;
    if (breakGlassReason) headers['x-break-glass-reason'] = breakGlassReason;
    // ... JSON body or FormData

    let response;
    try {
      response = await fetch(url, { method, headers, body: payload, credentials: 'same-origin' });
    } catch (networkError) {
      throw new ApiError(0, 'Cannot reach the MediChain server. Check that it is running.', null);
    }

    const expired = response.status === 401 && retry && !path.startsWith('/api/auth/');
    if (expired) {
      try {
        await refreshSession();
      } catch (refreshError) {
        accessToken = null;
        throw new ApiError(401, 'Your session has ended. Please log in again.', { code: 'SESSION_ENDED' });
      }
      return request(method, path, { body, query, form, retry: false });
    }
    return parse(response);
  }
\end{lstlisting}

\begin{lstlisting}[language=JavaScript,caption={One shared refresh (\file{frontend/js/api.js})}]
  function refreshSession() {
    if (!refreshing) {
      refreshing = fetch('/api/auth/refresh-token', { method: 'POST', credentials: 'same-origin' })
        .then(parse)
        .then((body) => {
          accessToken = body.data.token;
          return body.data.user;
        })
        .finally(() => { refreshing = null; });
    }
    return refreshing;
  }
\end{lstlisting}

\begin{howto}
\begin{itemize}
  \item \code{accessToken} is a variable inside the module's closure: never in
    \code{localStorage}, never reachable from other scripts.
  \item When any call gets a 401 (access token expired after 15 minutes), the
    client silently refreshes once and retries the call. The user never sees it.
  \item A dashboard page often makes five calls at once; if all five get 401,
    \code{refreshing} makes them share \emph{one} refresh request. This matters
    because refresh tokens rotate: five parallel refreshes would invalidate
    each other.
  \item A network failure becomes \code{ApiError} with status 0, which the
    screens show as ``Server unreachable''.
  \item While emergency access is switched on, every request carries the
    \code{x-break-glass-reason} header.
\end{itemize}
\end{howto}

\section{The dashboard shell: \file{js\codesl{}portal\codedot{}js}}

\subsection{Start-up}

\begin{lstlisting}[language=JavaScript,caption={Portal.start (\file{frontend/js/portal.js})}]
  async function start(config) {
    state.config = config;
    applyTheme(savedTheme());
    let user;
    try {
      user = await api.refreshSession();
    } catch (error) {
      if (error.status === 0) {
        document.getElementById('app').innerHTML = '<div class="boot-screen"><div class="state state-error"><h4>Server unreachable</h4><div>Start the MediChain API and reload this page.</div></div></div>';
        return;
      }
      sendToLogin();
      return;
    }
    if (user.role !== config.role) {
      location.replace(DASHBOARDS[user.role]);
      return;
    }
    state.user = user;
    document.getElementById('app').addEventListener('click', onClick); // one delegated listener for the whole app
    renderLayout();
    renderBanners();
    window.addEventListener('hashchange', () => showPage());
    await showPage();
    refreshBadges();
    state.pollTimer = setInterval(refreshBadges, 20000);
  }
\end{lstlisting}

Opening \file{doctor\codedot{}html} while logged in as a patient sends you to
\file{patient\codedot{}html}. This is a convenience only: the real protection is the
server's role check on every endpoint.

\subsection{Rendering a page: real loading, empty and error states}

\begin{lstlisting}[language=JavaScript,caption={showPage (\file{frontend/js/portal.js})}]
  async function showPage(pageId = currentPageId()) {
    const token = ++state.renderToken;
    state.page = pageId;
    state.handlers = {};
    state.charts.forEach((c) => c.destroy());
    state.charts = [];
    // ... highlight the nav item, set the title
    const container = document.getElementById('page-container');
    UI.loading(container);
    const render = state.config.pages[pageId] || SHARED_PAGES[pageId];
    try {
      await render(ctx(container, token));
    } catch (error) {
      if (token !== state.renderToken) return;
      if (error.code === 'SESSION_ENDED') return sendToLogin();
      UI.errorState(container, error, () => showPage(pageId));
    }
    if (token === state.renderToken) refreshIcons();
  }
\end{lstlisting}

\begin{howto}
The URL hash (\code{\#appointments}) selects the page. The shell shows a
loading state, then calls the role script's page function with a context
object \code{c}. If the function throws (any \code{ApiError}), the shell
replaces the page with an error state and a \emph{Try again} button. Nothing is
ever shown from a previous load or a fixture. The \code{renderToken} solves a
race: if the user clicks another page while the first is still loading, the
first page's late result is discarded (\code{c.stale()} becomes true).
\end{howto}

\begin{lstlisting}[language=JavaScript,caption={The error state (\file{frontend/js/ui.js})}]
  function errorState(container, error, onRetry) {
    const offline = error && error.status === 0;
    container.innerHTML = `<div class="state state-error" role="alert">
      ${icon(offline ? 'wifi-off' : 'alert-triangle')}
      <h4>${offline ? 'Server unreachable' : 'Could not load this page'}</h4>
      <div>${esc(error && error.message)}</div>
      ${onRetry ? '<button class="btn btn-secondary mt-3" data-retry>Try again</button>' : ''}
    </div>`;
    if (onRetry) container.querySelector('[data-retry]').addEventListener('click', onRetry);
    refreshIcons();
  }
\end{lstlisting}

\subsection{Escaping: the second line of XSS defence}

\begin{lstlisting}[language=JavaScript,caption={esc (\file{frontend/js/ui.js})}]
  function esc(value) {
    return String(value ?? '')
      .replace(/&/g, '&amp;').replace(/</g, '&lt;').replace(/>/g, '&gt;')
      .replace(/"/g, '&quot;').replace(/'/g, '&#39;');
  }
\end{lstlisting}

Pages build HTML with template strings, so \emph{every} value that came from the
server goes through \code{esc} before insertion. A patient named
\code{<img src=x onerror=alert(1)>} is displayed as text. A frontend test checks
this function, and the CSP (Section~\ref{ch:foundation}) blocks inline scripts
even if an escape were forgotten.

\section{A complete page: the patient's Permissions screen}

Every role page follows the same pattern: fetch in parallel, bail out if stale,
render, register click handlers.

\begin{lstlisting}[language=JavaScript,caption={Data loading and handlers (\file{frontend/js/patient.js}, abridged)}]
  async function permissions(c) {
    const [consents, doctors] = await Promise.all([c.api.get('/api/consents'), c.api.get('/api/doctors')]);
    if (c.stale()) return;
    const requested = consents.data.filter((g) => g.status === 'requested');
    const active = consents.data.filter((g) => g.status === 'active');
    // ... c.container.innerHTML = cards and tables built from these arrays,
    //     with buttons like <button data-action="revoke" data-id="...">

    c.on({
      approve: (el) => UI.modal({ /* choose a duration, then POST /api/consents/:id/approve */ }),
      deny: (el) => c.act(() => c.api.post(`/api/consents/${el.dataset.id}/deny`)),
      revoke: async (el) => {
        if (await UI.confirmDialog('Revoke access', 'The doctor will lose access to your records immediately.', { confirmLabel: 'Revoke', danger: true })) {
          await c.act(() => c.api.post(`/api/consents/${el.dataset.id}/revoke`));
        }
      },
      // ... grant
    });
  }
\end{lstlisting}

\begin{howto}
Buttons carry \code{data-action} and \code{data-id} attributes instead of
\code{onclick} code (which the CSP would block). The shell's single click
listener finds the nearest \code{[data-action]} element and calls the handler the
page registered with \code{c.on}. \code{c.act} runs an API call, shows a success
or error toast with the server's message, reloads the page and refreshes the
badge counters. Empty lists render through \code{UI.table(..., 'No doctor
currently has access')}, which shows a real empty state instead of an empty
table.
\end{howto}

\section{Login: \file{loginscript\codedot{}js}}

\begin{lstlisting}[language=JavaScript,caption={Two-step sign-in (\file{frontend/loginscript.js}, abridged)}]
  async function requestCode(button) {
    setLoading(button, true);
    try {
      const res = await api.post('/api/auth/login/request-otp', { email: pending.email, password: pending.password });
      $('verificationContact').textContent = res.data.email;
      $('otpMinutes').textContent = res.data.expiresInMinutes;
      return true;
    } catch (error) {
      const remaining = error.body && error.body.attemptsRemaining;
      setMessage('loginMessage', remaining !== undefined ? `${error.message}. ${remaining} attempt(s) left before the account locks.` : error.message);
      return false;
    } finally {
      setLoading(button, false);
    }
  }

  async function onVerify() {
    const otp = codeInputs().map((i) => i.value).join('');
    if (!/^\d{6}$/.test(otp)) { setMessage('otpMessage', 'Enter all 6 digits of the code.'); return; }
    // ...
      const res = await api.post('/api/auth/login/verify-otp', { email: pending.email, otp });
      pending.password = '';
      // ... redirect to the dashboard for res.data.user.role
  }

  function redirectAfterLogin(role) {
    const next = new URLSearchParams(location.search).get('next');
    // Only follow "next" when it points at this user's own dashboard (prevents open redirects).
    const target = next && next.startsWith(DASHBOARDS[role]) ? next : DASHBOARDS[role];
    location.replace(target);
  }
\end{lstlisting}

The password is kept in memory only between the two steps (so ``resend code''
works) and cleared as soon as the code is accepted. The six code boxes accept
typing, pasting a whole code, and backspace across boxes. The registration form
loads its insurer and department lists from \code{GET \codesl{}api\codesl{}settings\codesl{}public}.

\section{What each dashboard shows}

\begin{longtable}{@{}P{3.4cm}P{10.6cm}@{}}
\caption{Dashboard pages and their endpoints}\label{tab:pages}\\
\toprule
\textbf{Page} & \textbf{Content (endpoints)} \\
\midrule
\endfirsthead
\toprule
\textbf{Page} & \textbf{Content (endpoints)} \\
\midrule
\endhead
\multicolumn{2}{@{}l}{\textit{Patient}}\\
Dashboard & Next appointment, access requests, active medications, latest vitals, unpaid bills (appointments, prescriptions, consents, bills, vitals) \\
Appointments & Upcoming and history; book with a live slot grid; reschedule; cancel \\
Medical Records & Records with attachments, upload a document, download, ledger history, Verify integrity \\
Medications & Active and past prescriptions; request refill \\
Health Metrics & Charts of the last 7 days of vitals, abnormal-reading warning, registered devices, add a reading \\
Permissions & Pending requests, active grants, history; approve, deny, grant, revoke \\
Billing \& Claims & Bills, record payment, file claim, claim timeline \\
\multicolumn{2}{@{}l}{\textit{Doctor}}\\
Dashboard & Today's schedule, pending confirmations, refill queue, consent counts \\
Appointments & Today / Upcoming / Past; confirm, complete (with notes), no-show, cancel \\
My Patients & Patients with consent status; find by email; request access; patient chart with records, vitals, prescriptions, new record, upload, prescribe \\
Prescriptions & All prescriptions; approve or deny refills; complete or cancel \\
IoT Monitoring & Live tiles of each consenting patient's latest vitals, refreshed every 5 s \\
Emergency access & Header button that switches break-glass on with a reason \\
\multicolumn{2}{@{}l}{\textit{Admin}}\\
Dashboard / Reports & Live statistics and charts (\code{\codesl{}api\codesl{}admin\codesl{}stats}) \\
Verification & Pending doctors and agents; approve or reject \\
Users & Search, filter, activate/deactivate, create staff \\
IoT Devices & Registry; register (key shown once); activate/deactivate \\
Hospital Settings & Name, contacts, booking window, consent duration, insurers, departments \\
Ledger & Chain statistics, browse blocks, Verify ledger \\
Access Control & The access chain: allow/deny decisions, break-glass uses \\
\multicolumn{2}{@{}l}{\textit{Insurance agent}}\\
Dashboard & Queue size, amounts, approval rate \\
Claims Queue & Filter by status; process a claim through the state machine \\
Settlements, Reports & Settled claims; monthly and status aggregates with charts \\
\multicolumn{2}{@{}l}{\textit{All roles}}\\
Notifications & List, mark read, jump to the related page \\
Profile \& Security & Edit profile (and availability for doctors), change password, current trust score with its reasons \\
\bottomrule
\end{longtable}
